\section{A bounded search through alternative spanning-tree gauges}
\label{sec:cocycle-trees}

Section~\ref{sec:cocycle-connectivity} makes raw tree-gauged surfaces cheap
to certify. The remaining obstacle is discovery: a primitive surface and
its minimum-piece representative can both miss a disc. We now reuse one
cohomology solution to construct a bounded family of alternative tree
gauges. Five previously unsuccessful source records acquire native disc
proofs within the first three deterministic trials. An optional longer
search finds a sixth. The family remains incomplete.

\subsection{Changing the tree without solving cohomology again}

Fix an integral cocycle $c$ on the validated finite exterior. For a rooted
spanning tree $A$, define $p_A(v)$ by summing signed values of $c$ along
the unique tree path from the root to $v$. On every oriented tree edge
$u\to v$,
\[
 p_A(v)-p_A(u)=c(u,v).
\]
Consequently
\[
 c_A=c-\delta p_A
\]
vanishes on the tree and represents the same integral cohomology class.
If $h_{t,i}$ are the original local heights, the new heights are simply
$h_{t,i}-p_A(v_{t,i})$. Adding a common constant is immaterial. A different
tree can change the normal surface and its Euler characteristic, while
preserving the primitive class.

The independent raw certificate from the preceding section applies
unchanged: the zero-valued edges contain a spanning tree, so no nonempty
component union represents zero. Rank-one primitivity implies connectedness
and orientability. Euler characteristic one therefore supplies an essential
disc. The verifier reconstructs the signed cocycle, checks primitivity and
the connected zero-edge graph itself; it need not trust or reconstruct the
producer's choice of tree.

There is no need to repeat rational cohomology elimination. One verified
primitive cocycle supplies every gauge by path integration. Nor does each
tree introduce a different span-optimization problem: varying a subsequent
vertex potential still ranges over the same set
$\{c+\delta f\}$. The scheduler performs at most one span solve, on the
original heights. Alternative unoptimized representatives can be valuable
even when that common minimum-piece surface is not a disc.

\subsection{Deterministic discovery and bounded later exploration}

The producer recovers global oriented edge endpoints and values from every
local height difference, checks consistency, and builds the one-skeleton
adjacency lists. Loops cannot belong to a spanning tree and are omitted
from those lists; their cocycle values still enter the surface. Parallel
edges are retained. Roots are ordered by decreasing nonloop degree, with
global vertex label breaking ties. A breadth-first traversal chooses a
tree by the fixed adjacency order.

At most four such trials form the default prelude. Their root is the trial
index modulo the number of vertices. For explicitly requested trials after
the fourth, a private generator with seed 261009471 shuffles each adjacency
list before breadth-first traversal. This changes the candidate trees,
does not alter global random state, and supplies reproducible experiments
in the recorded Python environment. No probabilistic completeness claim
depends on these choices.

Normalize every integrated potential at the smallest global vertex.
Equality of the resulting tuples proves equality of the gauged cocycles,
so the producer skips duplicates, including the original zero potential.
The duplicate still consumes its trial allowance and construction work.
Each distinct candidate reconstructs binary normal coordinates and checks
normal matching before computing Euler characteristic. A positive always
passes the full independent source-bound verifier. A negative trial emits
no negative knot certificate; its connectedness follows from the previously
checked primitive source class and the tree construction.

With $V,E=O(T)$ vertices and edges, graph setup and root ordering use
$O(T\log(T+1))$ comparison work after manifold validation. Each of the
first four deterministic trials uses $O(T)$ arithmetic and incidence work.
If the original heights have $B$ bits, a tree path has at most $V-1$ edges,
so potentials and new heights have $O(B+\log(T+1))$ bits. Binary coordinate
construction and the Euler calculation do not expand normal pieces.
Storing the graph and one candidate takes $O(T)$ records; retaining exact
duplicate signatures for $R$ requested trials takes $O(RV)$ integers.
The optional shuffled trials add the permutation-generation cost of
Python's random shuffle; the $O(T)$ traversal bound is conditional on those
adjacency orders already being chosen. No worst-case bound for the
optional random generator is being used to support a recognition theorem.

\subsection{Switching between search methods}

The optional native stage now follows this bounded schedule:
\begin{enumerate}
\item Check the initial primitive cocycle. A success ends the stage
      before constructing the tree-search graph.
\item Try up to four requested tree gauges, skipping duplicate potentials
      and stopping immediately on a verified disc.
\item If enabled, minimize the original span and check its surface.
\item On another miss, try any remaining requested tree gauges.
\end{enumerate}
All steps share the existing cooperative work allowance and the enclosing
deadline. A cap or unsuccessful search returns \code{INCONCLUSIVE} to the
exact fallback. Thus cheap early successes avoid flow search, and a
successful flow avoids every optional later tree. The schedule uses the
outcome of completed attempts; it does not infer that a small piece count
or improved Euler characteristic predicts future success.

The API options are \code{tree\_trials} for \code{normal\_seed\_decide}
and \code{normal\_seed\_tree\_trials} for \code{recognize}; the CLI uses
\code{--normal-seed-tree-trials}. The default is four. Zero restores the
previous raw-then-span search. \code{--normal-seed-no-optimize} now skips
the span solve; set the tree allowance to zero as well to request only
the original surface. The stage as a whole remains disabled by default.
Boolean, negative and noninteger trial allowances are rejected even when
the stage is disabled.

\subsection{Exploratory alternatives and their limits}

Before integration, a standalone experiment tried up to 24 tree gauges on
53 previously unsuccessful raw candidates selected from the existing
source corpus. It found six discs; five appeared in deterministic trials
two or three, and one in shuffled trial 22. These observations motivated
the short default prelude and optional later search. The maintained audit
below repeats all 84 source records, checks old-proof preservation and
validates the new surfaces independently. Euler changes in the exploratory
log are not treated as genus changes: boundary-component counts can vary.

A second experiment minimized the sum of absolute global-edge cocycle
values. Each term $|c(e)+f(v)-f(u)|$ can be represented as a two-height
span, embedded in the existing four-corner arithmetic API by repeating
both endpoints. This objective also strictly decreases on deletion of a
nonempty zero-class component union, since every normal piece meets edges.
Its rank-one primitive optima therefore have the same connectedness
argument. But the 58 tested old-stage misses of at most twelve crossings
produced no Euler-one surface under the two-million-guard optimization
allowance; all calls completed. The experiment is retained, without adding
another optimization stage or source-proof format to the recognizer.

Neither alternative searches every normal surface. In particular, the
three-crossing unknot $(1,1,-1)$ remains a miss under the default policy.
There is still no certified complete search with a general subexponential
or quasipolynomial running-time bound.

\subsection{New source proofs and preserved controls}

The 22 focused tree, seed and connectivity tests pass in 3.227 seconds.
The full suite passes all 1,187 tests in 229.784 seconds. Tests establish
that early successes avoid flow, raw successes avoid tree setup, optimized
successes avoid later trials, and the late example appears after the
optimization miss. They check exact option types, duplicate skipping,
determinism, matching failures, shared-cap exhaustion and cancellation.
Positive replay succeeds with tree discovery disabled and uses the
unchanged independent verifier.

All 84 source records are run with zero, four and 24 tree trials. The zero
policy reproduces the preceding stage's 18 positive records, 65 completed
restricted misses and one cap. Four trials produce 23 positives and 60
completed misses; 24 trials produce 24 positives and 59 completed misses.
Gordian remains capped in both policies. All eighteen old positive hashes
are preserved, and all six new source witnesses replay independently.
The following rows list the new successes; trial counts include duplicates.

\begin{center}
\small
\begin{tabular}{lrrrrr}
\toprule
Source & $n$ & Trial & Disc pieces & Minimum pieces & Guards \\
\midrule
random-29 & 4 & 3 & 83 & 63 & 11,227 \\
random-79 & 4 & 2 & 102 & 63 & 10,192 \\
random-99 & 3 & 2 & 66 & 43 & 7,785 \\
random-124 & 5 & 3 & 121 & 79 & 14,185 \\
random-150 & 7 & 2 & 153 & 99 & 18,281 \\
random-81 (24 trials) & 5 & 22 & 122 & 71 & 180,578 \\
\bottomrule
\end{tabular}
\end{center}

Every displayed disc has more pieces than the span optimizer's returned
non-disc. For example, \code{random-29} changes from an 89-piece raw
surface of Euler characteristic zero to an 83-piece disc in trial three;
the checked minimum-piece solution has 63 pieces and Euler characteristic
$-1$. This proves that a useful disc exists within the same potential
family even though the selected optimum is not a disc. It does not prove
that every tied minimum-piece solution is a non-disc.

Fresh Regina controls check all 205 distinct generated candidates from
the six new source examples and three further controls: trefoil,
figure-eight and the genus-one unknot miss. Independent native inspection
and Regina agree on connectedness, orientability, Euler characteristic
and essential-disc count in every case. The new source certificates do
not depend on Regina. The full audit takes 84.487 seconds while the
test suite also runs, so that elapsed time is not a performance comparison.

Nineteen ordinary corpus inputs give their expected verdicts with zero,
four or 24 trials. All 57 calls finish before reaching the native stage.
Thus the additional source coverage is a property of the optional native
search, not a demonstrated expansion of the ordinary recognizer's solved
corpus or a general running-time improvement.

\subsection{Complete recognition gains and unsuccessful-search costs}

The benchmark compares the published raw-then-span stage with the new
four-trial default, except for the explicitly labeled 24-trial case. Each
timed endpoint parses the input PD, runs the complete recognizer and exact
fallback when needed, independently replays any returned native proof,
and hashes and counts its serialized certificate. Earlier shortcuts are
disabled to make the native stage reachable. Both arms use no local seed
cap and retain a twelve-second global deadline. The native stage's
ordinary disabled default is not being changed by this benchmark.

Runs are serial, without overlapping tests, audits, builds or commits.
Every round randomly orders old, old A/A, new and new A/A calls. The initial
five-round experiment completes all 160 measured and 32 warm-up calls in
16.851 seconds, but its unchanged raw-success control has a new-arm A/A
ratio of 1.237 and an old/new ratio of 0.863. That variation warrants a
repeat before interpreting the control. The full initial record remains
in \path{data/cocycle-trees-initial-recognize.json}.

The final nine-round experiment completes all 288 measured and 32 warm-up
calls in 28.078 seconds. The table gives median primary-arm times and
median paired ratios; these ratios need not equal ratios of median times.
All verdicts and each implementation's proof hashes agree across repeated
calls and both experiments. Existing positive proofs remain byte-identical
across implementations. Newly successful native calls intentionally replace
the old inconclusive stage followed by an exact fallback.

\begin{center}
\small\begin{tabular}{lrrrrr}
\toprule
Input & Old (ms) & New (ms) & Old/new & Old A/A & New A/A \\
\midrule
Raw circle & 65.061 & 63.657 & 1.026 & 1.025 & 1.005 \\
Optimized positive & 129.131 & 144.014 & 0.925 & 0.991 & 0.978 \\
Early trial 2 & 50.894 & 38.314 & 1.345 & 1.017 & 0.988 \\
Early trial 3 & 78.248 & 51.757 & 1.516 & 0.995 & 0.991 \\
Five-crossing early & 110.901 & 65.990 & 1.639 & 0.958 & 0.989 \\
Late, 24-trial option & 111.954 & 181.031 & 0.624 & 1.024 & 1.014 \\
Genus-one miss & 51.600 & 59.570 & 0.866 & 1.004 & 1.000 \\
Trefoil & 50.198 & 57.041 & 0.877 & 0.982 & 1.024 \\
\bottomrule
\end{tabular}

\end{center}

The three early successes improve complete recognition by
$1.35$--$1.64\times$. They are \code{random-99}, \code{random-29} and
\code{random-124} from the preceding table. The unchanged raw-success
control is near parity: its 1.026 paired ratio is comparable to its 1.025
old-arm A/A ratio. There is no measured algorithmic gain to claim for
that branch.

The unsuccessful prelude has a visible cost. The previous optimized
positive is about eight percent slower by the paired ratio, while the
genus-one miss and trefoil cost about 15 and 14 percent more respectively.
Both misses still use the exact Khovanov fallback and return their expected
\code{UNKNOT} and \code{KNOTTED} verdicts. The optional 24-trial example
produces a new native disc proof, but its full call is about 60 percent
slower than the old stage plus fallback. The initial run also shows this
cost, with a paired ratio of 0.540. This is additional native proof
coverage, not a speedup for that example; the longer search is therefore
an explicit option rather than the default allowance.

These measurements support a bounded cheap prelude with a documented
tradeoff. They do not establish a uniform speedup, a solved-corpus gain for
the ordinary portfolio, or a global complexity improvement. The zero-trial
option retains the earlier candidate schedule when its cost is preferable.


Reproduction uses \path{../fast/normal_orbit_research/trees.py}, with
\code{audit} and \code{recognize} modes; the final timing run requests
\code{--rounds 9}. The baseline is \code{48ef0ccef}. All 481 runtime,
driver, fixture and prior-audit pins agree. The earlier exploratory scripts
and full logs are retained in \path{data/tree-gauge-exploration.py} and
\path{data/edge-weight-exploration.py}, with corresponding text outputs.
Run those scripts from the \code{fast} directory with its Python package
on \code{PYTHONPATH}. The incoming archives already have maintained review
records, and the checked \code{openai/math} head remains \code{fd4aeeb2e}.

